\section{Conclusions and Future Work}
% =============================================================

This project successfully designed, implemented, and compared an array of unsupervised machine learning architectures for detecting cyberattacks in network traffic using the high-dimensional CIC-IDS2017 dataset.

The main conclusions of our investigation are:
\begin{itemize}
  \item Classical unsupervised anomaly detection methods (LOF, OC-SVM, Isolation Forest) struggle severely in 70-dimensional spaces, validating the curse of dimensionality.
  \item PyTorch-based Deep Autoencoders overcome this bottleneck by mapping features into low-dimensional latent spaces, learning robust non-linear boundaries.
  \item The proposed **Ensemble Autoencoder** achieves state-of-the-art anomaly discrimination (AUC = $0.9810$) and successfully catches **$95.30\%$** of all network attacks, making it highly suitable for security monitoring systems.
  \item The **Second Autoencoder** (anomaly-trained) is the most operationally balanced model, achieving an F1 of $0.7844$ and a precision of $77.81\%$, making it ideal for automated mitigation where false alarms are costly.
\end{itemize}

Future work will focus on:
\begin{enumerate}
  \item \textbf{Hybrid Architectures:} Exploring Deep Autoencoding Gaussian Mixture Models (DAGMM) to combine reconstruction error with density energy.
  \item \textbf{Attention Mechanisms:} Incorporating attention layers in the encoder to let the model focus on key temporal traffic features.
  \item \textbf{Weighted Ensembles:} Implementing weighted scoring logic in the ensemble instead of simple binary OR logic to improve precision while maintaining high recall.
\end{enumerate}
